\begin{figure}[htbp]
\centering
\begin{tikzpicture}[beach diagram,x=0.43cm,y=0.45cm]
  \foreach \offset/\ratio/\r/\title in {0/0.4/0.6/単調に減衰,12/1.4/-0.4/交互に減衰,24/2.2/-1.2/交互に増大} {
    \begin{scope}[xshift=\offset*0.43cm]
      \node[beach heading] at (4,6.2) {\title};
      \node at (4,5.1) {$h/\tau=\ratio$};
      \draw[black!18] (0,-4.5) rectangle (8,4.5);
      \draw[beach axis] (0,0)--(8.5,0) node[right] {$n$};
      \draw[beach axis] (0,-4.5)--(0,4.7);
      \foreach \y in {-4,-2,2,4} {
        \draw (0,\y)--(-0.12,\y) node[left] {$\y$};
      }
      \foreach \n in {2,4,6,8} {
        \draw (\n,0)--(\n,-0.14);
        \node[below] at (\n,-4.6) {$\n$};
      }
      \foreach \n in {0,...,7} {
        \pgfmathsetmacro{\ya}{pow(\r,\n)}
        \pgfmathsetmacro{\yb}{pow(\r,\n+1)}
        \draw[BeachBlue,thick] (\n,\ya)--(\n+1,\yb);
      }
      \foreach \n in {0,...,8} {
        \pgfmathsetmacro{\yn}{pow(\r,\n)}
        \fill[BeachBlue] (\n,\yn) circle[radius=1.4pt];
      }
    \end{scope}
  }
  \node[rotate=90] at (-2.2,0) {固定点からのずれ $e_n/e_0$};
  \node[beach note] at (16,-6.2) {$e_{n+1}=(1-h/\tau)e_n$、$e_n/e_0=(1-h/\tau)^n$};
\end{tikzpicture}
\caption[単一線形緩和モデルのバッチ幅と安定性]{実負の固有値 $-1/\tau$ を一つだけ持つ線形モデルの、バッチ幅による安定性の違い。
各点は図中の漸化式から直接計算した値で、BEACH の帯電履歴ではない。
全 panel の縦軸尺度は共通である。標本誤差、複素固有値、非正規結合は含めていない。}
\label{fig:batch_stability}
\end{figure}
